\scanpage{62}
b) L'application $\Sigma^{<}(2)\xrightarrow{(p_1\circ s',\unsure{p_1b'})}\Sigma\times\Sigma$
se factorise par $\Sigma^{<}(1)\lhook\joinrel\xrightarrow{(p_1,\uline{b})}\Sigma\times\Sigma$
en un morphisme
\[
p_2:\Sigma^{<}(2)\longrightarrow\Sigma^{<}(1)
\]
qui rende le carré
\diagram{\begin{tikzcd}[column sep=small,execute at end picture={
\draw[densely dotted] ([xshift=7pt]s62sourcea.east) arc[start angle=270,end angle=90,radius=3pt];
\draw[densely dotted] ([xshift=7pt]s62sourceb.east) arc[start angle=270,end angle=90,radius=3pt];}]
|[alias=s62sourcea]|\Sigma^{<}(2)\arrow[rr,"\uline{b'}",start anchor=east,shorten <=7pt]\arrow[dr,"p_2"'] && \Sigma^{<}(1)\arrow[dr,"p_1"]&\\
& |[alias=s62sourceb]|\Sigma^{<}(1)\arrow[rr,"\uline{b}"',start anchor=east,shorten <=7pt] && \Sigma
\end{tikzcd}}
\uline{Cartésien}, et $p_1:\Sigma^{<}(1)\longrightarrow\Sigma$ étant transversal
à \deleted{\unsure{l'imm}}ersion-but $\uline{b}:\Sigma^{<}(1)\longrightarrow\Sigma$.
\margin{\uline{NB} (Si $\Sigma$ est séparé, donc $\Sigma\times\Sigma$ séparé
sur $\Sigma$, et $\Sigma^{<}(1)\xrightarrow{b}\Sigma$
propre, il en résulte
que $\Sigma^{<}(1)\longrightarrow\Sigma\times\Sigma$ est
propre donc un plongement
\uline{fermé}.}
\margin{cette condition de transversalité est
peut-être conséquence des autres\ldots}

De façon équivalente, une telle structure
est donnée par un espace topologique \inserted{séparé} $\Sigma$,
muni d'une relation d'ordre $x\preceq y$
\deleted{dont le graphe est}\inserted{satisfaisant les conditions suivantes}

\deleted{a) le graphe \deleted{$\unsure{\Sigma^{<}(1)}$} $\Sigma(1)$ de cette relation est
fermé dans $\Sigma\times\Sigma$, et $\Sigma(1)\xrightarrow{b}\Sigma$ est
une immersion.
\deleted{$\unsure{\Sigma}$}
b) Si $\Sigma^{*}(1)$ est le graphe de la
relation d'ordre \inserted{stricte}, $\Sigma^{*}(1)\xrightarrow{\uline{b}}\Sigma$ est une
immersion \inserted{bordante} propre \deleted{\unsure{(donc la partie de}}
a) Si $x\in\Sigma$, $\Sigma_{<x}$ \inserted{$(=\Sigma^{*}(1)_x)$}, \deleted{\unsure{est}}
est (fini) \uline{tot${}^t$.} \uline{ordonné}.
}
\mspage{61}
